% ======================================================
%  PART II — CÁC GIẢI PHÁP THỰC THỂ (Conceptual Solutions)
% ======================================================
% Yêu cầu: Với mỗi problem ở Part I, đề xuất 2-3 conceptual solutions
% KHÁC BIỆT ĐÁNG KỂ so với nhau.
%
% Không phải 2-3 version sửa đổi nhỏ — mà là các hướng tiếp cận
% khác nhau về fundamental approach.
\section{Conceptual Solutions}

\subsection{Vấn đề 1: Trì hoãn (\textit{lag}) khi xem trực tiếp}

\subsubsection{Concept A: Sử dụng trí tuệ nhân tạo để phân đoạn nội dung phát trực tiếp}

Trí tuệ nhân tạo có thể dựa vào lời nói của bình luận viên và hình ảnh để phân đoạn nội dung phát trực tiếp thành các phần quan trọng và không quan trọng. 

Những đoạn quan trọng sẽ có ``dấu trang'' để người dùng nhảy đến ngay để xem lại. Dấu trang sẽ được làm nổi bật để tăng tỉ lệ người dùng nhấp vào. Khi đó, tỉ lệ click vào các đoạn không quan trọng giảm theo, giúp giảm tải không cần thiết cho hệ thống (không cần load những phần kém quan trọng).

Có ``fast mode'' để những đoạn không quan trọng chỉ hiện một khung hình tĩnh kèm text mô tả và giọng đọc của bình luận viên. Text sẽ được hiện dưới dạng streaming để người dùng có cảm giác rằng họ vẫn đang được cập nhật thông tin theo thời gian thực.

Nhược điểm là tốn chi phí AI (dù không nhiều vì không phải lần nào truyền dữ liệu video đến người dùng cũng phải phân tích AI, chỉ cần phân tích một lần và lưu lại). Với lại những người dùng đầu tiên join vào có thể sẽ chưa phân tích kịp nên fast mode chưa available.

% TODO:
% - Mô tả giải pháp ở mức khái niệm
% - Giải pháp này giải quyết được vấn đề như thế nào?
% - Nhược điểm chính

\subsubsection{Concept B: Biến mỗi người dùng thành một kênh phát trực tiếp riêng}
% TODO: Mô tả — PHẢI khác đáng kể so với Concept A
% Ví dụ: Concept A là "tự động", Concept B là "user-controlled"

Ý tưởng là ứng dụng sẽ tạo chức năng video call cơ bản dạng P2P -- kết nối giữa người dùng với nhau không thông qua server chính. Khi những người dùng sớm nhất join vào xem một giải đấu lớn, ứng dụng sẽ hỏi là người dùng có muốn bật tính năng video call -- cho phép người dùng khác xem cùng thông qua video người dùng hiện tại đang chia sẻ trong call hay không, với đặc quyền được ưu tiên về băng thông tốc độ cao (trên UI nhớ nhấn mạnh cái ý này để người ta đồng ý). Khi đó, những người dùng sau sẽ có tùy chọn join vào xem thông qua video call của người dùng đầu tiên, giúp giảm tải cho server chính.

Nhược điểm là giải pháp này chỉ khả thi khi số lượng người dùng sớm join vào đủ lớn để tạo ra một mạng lưới video call. Nếu không, những người dùng sau sẽ vẫn phải xem trực tiếp từ server chính, dẫn đến tình trạng lag vẫn xảy ra. Với lại những người xem chung có thể sẽ nghe phải âm thanh của những người dùng khác join cùng một call, nên cần có cơ chế tắt tiếng những người dùng khác.

% \subsection{Vấn đề 2: Giao diện quá nhiều hình ảnh, thiếu thông tin văn bản}
% % TODO: Tương tự — 2-3 concepts khác nhau
%
% \subsubsection{Concept A: \textit{[Tên giải pháp]}}
% % TODO
%
% \subsubsection{Concept B: \textit{[Tên giải pháp]}}
% % TODO

\subsection{Vấn đề 2: Quá tải thông tin đối với người dùng mới}
% TODO: Tương tự — 2-3 concepts khác nhau

Có các vấn đề: Nhiều banner và carousel quá nhưng thực chất có thể không đúng cái mà người dùng muốn xem, đặc biệt với người dùng lâu lâu mới xem và fan theo mùa -- biết mình cần xem gì ngay từ đầu rồi; Quá nhiều category trong cùng một trang.

\subsubsection{Concept A: Rút số lượng banner và carousel}
% TODO
Từ việc sử dụng nhiều banner và carousel, nhóm chuyển sang dùng 1 carousel duy nhất và đặt bên sidebar ở trái hoặc phải để thứ nhất là vẫn không bị chán, có animation chuyển qua lại giữa các slide mà vẫn giảm mật độ nội dung. Điều này dựa trên insight rằng người dùng được khảo sát thường bỏ qua carousel vì tính cá nhân hóa không cao (họ là người dùng mới). Có khả năng những người dùng mới khác cũng sẽ giống như họ. 

Nhược điểm là bị chiếm mất chiều rộng của màn hình, nhưng có thể bù đắp bằng cách cho phép thu gọn carousel dạng dọc này lại.

\subsubsection{Concept B: Cho phép người dùng vào edit mode để drag and drop các thành phần trên giao diện}
% TODO
Người dùng có thể vào edit mode để kéo thả vào trang chủ những nội dung, chủ đề phim mình muốn hiển thị. Có thể ẩn đi những thứ mình không muốn như banner, carousel, thumbnail của một thể loại phim nào đó, một môn thể thao nào đó. Khi đó, người dùng mới sẽ không bị quá tải thông tin vì họ có thể tự sắp xếp giao diện theo ý mình. Điểm khác biệt so với filter thuần là có thể thay đổi thứ tự của các thành phần trên giao diện, không chỉ ẩn hiện. Với lại sau khi save thì lần nào mở lên cũng sẽ có layout y chang như vậy chứ không bị reset.

Nhược điểm là tùy chọn này sẽ phù hợp hơn với người có kinh nghiệm chút ít về máy tính hoặc người xem thường xuyên, vì nếu không xem thường xuyên thì họ mất thời gian chỉnh giao diện rồi lưu lại làm gì khi mấy tháng mới vô xem một lần???